% ============================================================================
\section{The Temple wearable: technical analysis}
\label{sec:temple}

The project started from a request to gather ``every single technical detail'' of Temple, both hardware and software, by aggressive web research and engineering reasoning. Facts were tagged \emph{confirmed} (Temple's own preprints or website), \emph{reported} (press, sometimes conflicting) or \emph{inferred} (our engineering estimate).

\subsection{Company and research setup}

Temple is not a Zomato or Eternal product. It began in 2025 as a project of Goyal's personal lab, Continue Research ($\approx$\$25M of his own money), which also published the ``Gravity Ageing Hypothesis''. Goyal stepped down as Eternal CEO in January 2026. In February 2026 Temple raised \$54M at a $\approx$\$190M valuation (Goyal, Steadview, Peak XV, Info Edge Ventures, Dharana; angels including Kunal Shah, Vijay Shekhar Sharma, Nithin and Nikhil Kamath); a $\approx$\$500M valuation was reportedly targeted in July 2026. On 27 August 2026 it acquired the London longevity clinic Longevous (Dr Robert Mohr as CMO, Dr Avi Roy as Head of Science) to run trials. Entities are Temple Pvt.\ Ltd.\ (Gurugram) and Temple Wearables Inc.\ (San Francisco). David A.\ Boas (diffuse optics) sits on its Scientific Advisory Board; Satchidananda Panda (Salk) is a paid consultant.

\subsection{Hardware}

\begin{table}[H]
\centering\small
\caption{Temple hardware (confidence in brackets).}
\begin{tabularx}{\linewidth}{@{}l L@{}}
\toprule
Item & Detail \\
\midrule
Sensing principle & Multi-wavelength reflectance PPG [confirmed] \\
Wavelengths & Green (shallow vessels), red and infrared (deeper); exact nm not disclosed [confirmed] \\
Placement & Front of the temple over the superficial temporal artery, a scalp artery that does not feed the brain [confirmed] \\
Attachment & Medical-grade double-sided adhesive, lifetime supply [confirmed/reported] \\
Raw PPG sampling & 32\,Hz, green channel, right temple (age study) [confirmed] \\
Outputs & Brain Flow index 1\,Hz; Entropy every second [confirmed] \\
Heart-rate error & Median 3.6--3.8\% during exercise and posture changes [confirmed] \\
Geometry & Source--detector spacing withheld (listed as a limitation) [confirmed] \\
Form factor & Pebble-shaped metallic pod; v2 under half v1's volume [reported] \\
Battery, chipset, radio, IP rating & Not published \\
\bottomrule
\end{tabularx}
\end{table}

A January 2026 exchange with neurologist Dr Sudhir Kumar questioned the placement (the internal carotid at the neck does supply the brain); Goyal answered that PPG resolution is ``significantly higher at the temple'' and that ML models trained against MRI and transcranial Doppler bridge the gap.

\subsection{Software and metrics}

The \emph{Brain Flow index} (BFI) is unitless, relative to the user's baseline, and combines PPG perfusion, multi-wavelength haemoglobin changes and pulse-waveform shape, reportedly calibrated against TCD and MRI. \emph{Entropy} (June 2026) is a 1--250 score for ``metabolic and sympathetic demand'' with Maxima/Minima variants; a claimed r = 0.93 against a metabolic cart (heart rate alone r = 0.55) is unpublished. The age-study pipeline (confirmed) filters green PPG 0.5--10\,Hz (4th-order zero-phase Butterworth), resamples to 75\,Hz, cuts 25\,s windows, extracts 89 pulse-shape features (timing, amplitude, area, derivative ratios such as $(A_d-A_b)/A_a$ and AGImod, harmonic ratios) and reduces them by PCA (five components explain 85.3\%).

\subsection{Validation preprints}

\begin{table}[H]
\centering\small
\caption{Temple's three funded preprints (none peer-reviewed).}
\begin{tabularx}{\linewidth}{@{}L L L@{}}
\toprule
Study & Setup & Key result \\
\midrule
vs time-domain NIRS (bioRxiv, Sep 2026) & 44 young men, 60 sessions; tilt, squat, lying & Brain vs scalp layer correlation: tilt 0.91 vs 0.49; squat 0.84 vs 0.70; lying 0.84 vs 0.77; direction matched in 59/60 \\
vs transcranial Doppler (medRxiv, Sep 2026) & 23 adults 18--30 (20 per protocol); cycling, posture & Same-time r $\approx$0.80; lag-optimised 0.84 (exercise, $-$19\,s) and 0.82 (posture, +9\,s) \\
vs wrist PPG as an age signal (bioRxiv, Aug 2026) & $\approx$124 temple vs 240 wrist users & 14 pulse-shape features tracked age at the temple vs 3 at the wrist (different people at each site) \\
\bottomrule
\end{tabularx}
\end{table}

The authors themselves note that the signal is extracranial, MCA velocity is not total flow, only healthy young adults were tested, the algorithm is closed and the studies are industry-funded. The product was expected at $\approx$\$1,000 in Q4 2026 with no FDA or CDSCO clearance, positioned as general wellness; job postings mentioned BCI (EEG/EMG) engineering, but no EEG hardware is confirmed.

\subsection{Engineering inference and relevance to epilepsy}

We estimated (unconfirmed) a multi-LED/multi-photodiode optical front end (MAX86141/ADPD4100 class), an nRF52/nRF54-class BLE processor, a 6-axis IMU, a 20--50\,mAh cell (1--3 days), on-device filtering with BFI/Entropy computed in the app or cloud.

\textbf{Conclusion: Temple cannot detect seizures as shipped.} It measures optical blood-flow and pulse signals over a scalp artery, not the brain's electrical activity; at most it could pick up peri-ictal heart-rate and blood-flow changes, which nobody has validated. A device that combines optical sensing with EEG electrodes would fill a real gap; this redirected the project towards EEG plus body signals and, later, towards the ear (Section~\ref{sec:device}).
